\documentclass[11pt]{scrartcl}
\usepackage[a4paper,margin=2.6cm]{geometry}
\usepackage{../common/rlstudy}

\title{Extended reading --- Week 6: Rainbow DQN and BBF}
\author{STUDY-RL}
\date{\today}
\begin{document}
\maketitle

\section*{Reading list}

\subsection*{1. Hessel et al.\ 2018 --- ``Rainbow: Combining Improvements in Deep RL''}

\textbf{Citation.} Matteo Hessel, Joseph Modayil, Hado van Hasselt, Tom Schaul,
Georg Ostrovski, Will Dabney, Dan Horgan, Bilal Piot, Mohammad Azar, David Silver,
\emph{AAAI} 2018, arXiv:1710.02298.

\paragraph{Summary.} Combines Double + Dueling + PER + n-step ($n=3$) + NoisyNets + C51
into a single agent. Reports a careful 1-out ablation across 57 Atari games. Key finding:
multi-step ($n$) and the distributional value head are the two biggest contributors;
PER is third; NoisyNets, Dueling, and Double are roughly tied at a smaller magnitude.

\paragraph{Why it matters for this week.} The agent we implement (\texttt{project/agent.py})
is a deliberate Rainbow-lite: we keep Double, Dueling, PER, NoisyNets but skip the
distributional head and n-step (deferred to the workbook math).

\begin{quote}\itshape
``Our results show that combining the studied improvements yields the state-of-the-art on
the Atari 2600 benchmark, both in terms of data efficiency and final performance, while
giving up no critical individual contribution.''
\end{quote}
\textbf{Commentary.} The right way to read Rainbow is not as ``the algorithm'' but as a
\emph{template for combining orthogonal patches}. Each component fixes a different
failure mode; they compose because the failures are independent.

\subsection*{2. van Hasselt, Guez, Silver 2016 --- ``Deep Reinforcement Learning with Double Q-learning''}

\textbf{Citation.} Hado van Hasselt, Arthur Guez, David Silver, \emph{AAAI} 2016,
arXiv:1509.06461.

\paragraph{Summary.} Lifts the Double-Q idea from Hasselt 2010 (Week 4) into the
function-approximation regime. Concretely: use the target network for evaluation but the
online network for the argmax. Demonstrates that vanilla DQN over-estimates Q-values
across the entire Atari suite (running average of max-Q diverges over a few million
steps); Double-DQN keeps them bounded and improves final score on 47/57 games.

\paragraph{Why it matters for this week.} The single most impactful one-line patch on
DQN, and the template for the same trick later applied to actor-critics (TD3, REDQ).

\begin{quote}\itshape
``We show that this idea, which was originally proposed in a tabular setting, can be
generalized to work with large-scale function approximation. We propose a specific
adaptation to the DQN algorithm and show that the resulting algorithm not only reduces
the observed overestimations, as hypothesized, but that this also leads to much better
performance on several games.''
\end{quote}
\textbf{Commentary.} Note the structural parallel to van Seijen et al.\ 2009 (Expected
SARSA): the same decoupling-selection-from-evaluation trick reappears as soon as you
look at the bootstrap target through Jensen's lens.

\subsection*{3. Schaul et al.\ 2016 --- ``Prioritized Experience Replay''}

\textbf{Citation.} Tom Schaul, John Quan, Ioannis Antonoglou, David Silver, ICLR 2016,
arXiv:1511.05952.

\paragraph{Summary.} Replaces uniform replay sampling with priority $p_i \propto
|\delta_i|^\alpha$, importance-weight-corrected with $w_i = (N P(i))^{-\beta}$,
$\beta: 0.4 \to 1.0$. Introduces a sumtree for $O(\log N)$ sampling. Big improvements on
sparse-reward Atari games (Frostbite, Montezuma) where uniform sampling rarely sees the
informative transition.

\paragraph{Why it matters for this week.} The IS-weight derivation in our notes follows
this paper directly, and the algorithm is implemented in \texttt{common/replay.py} as
\texttt{PrioritizedReplayBuffer}.

\begin{quote}\itshape
``With prioritized replay we develop a framework for prioritizing experience, so as to
replay important transitions more frequently, and therefore learn more efficiently. We
use prioritized experience replay in Deep Q-Networks (DQN), a reinforcement learning
algorithm that achieved human-level performance across many Atari games.''
\end{quote}
\textbf{Commentary.} Note how the importance-sampling story connects to OPE. The
mechanism is identical: weight each sample by the ratio of the target distribution to
the sampling distribution. In PER the target is uniform over the buffer; in OPE the
target is the policy we want to evaluate.

\subsection*{4. Bellemare, Dabney, Munos 2017 --- ``A Distributional Perspective on RL''}

\textbf{Citation.} Marc G.\ Bellemare, Will Dabney, R\'emi Munos, \emph{ICML} 2017,
arXiv:1707.06887.

\paragraph{Summary.} Introduces C51 (categorical 51-atom distributional DQN). Models the
return distribution $Z(s,a)$ rather than just its mean $Q(s,a)$. Trains with a
KL-projected distributional Bellman backup; chooses actions by $\argmax \E_Z[Z]$ but
benefits from the richer auxiliary signal (better representations, more stable targets
after the $\min$/$\max$ aggregation). Beats DQN on 41/57 Atari games at $V_{\min}=-10$,
$V_{\max}=+10$, 51 atoms.

\paragraph{Why it matters for this week.} The conceptual jump from ``regression on a
mean'' to ``classification on a distribution''. The 2024 Farebrother ``Stop Regressing''
paper (HL-Gauss two-hot loss) is the direct descendant in continuous control.

\begin{quote}\itshape
``We argue for the fundamental importance of the value distribution: the distribution of
the random return received by a reinforcement learning agent. This is in contrast to the
common approach to reinforcement learning that models the expectation of this return, or
value.''
\end{quote}
\textbf{Commentary.} The framing is more important than the specific algorithm. Once you
internalise that the right learning target is a distribution, risk-sensitive control,
uncertainty quantification, and exploration bonuses all become natural; with a scalar
mean they require auxiliary machinery. See \texttt{research/deeprl\_core\_sota.md} \S4
for the modern landscape.

\subsection*{5. Schwarzer et al.\ 2023 --- ``Bigger, Better, Faster: Human-Level Atari with Human-Level Efficiency'' (BBF)}

\textbf{Citation.} Max Schwarzer, Johan Obando-Ceron, Aaron Courville, Marc Bellemare,
Rishabh Agarwal, Pablo Samuel Castro, \emph{ICML} 2023, arXiv:2305.19452.

\paragraph{Summary.} Combines Rainbow with three modern ingredients: (i) a wider
ResNet backbone (15-layer Impala-style with width multiplier), (ii) periodic
\emph{shrink-and-perturb resets} every 40k gradient steps,
(iii) discount and n-step \emph{annealing} ($\Disc: 0.97 \to 0.997$, $n: 10 \to 3$). EMA
target update at $\tau=0.005$. Self-prediction (SPR) auxiliary loss. Reaches super-human
median on Atari-100k in around 6 GPU-hours --- the current single-GPU SOTA for
sample-efficient discrete control.

\paragraph{Why it matters for this week.} The 2025 frontier. Read this paper as the
``what is Rainbow now?'' answer.

\begin{quote}\itshape
``We dissect a long-standing puzzle in deep RL: why don't bigger networks help? We argue
that the answer is plasticity loss, which we address with periodic shrink-and-perturb
resets. With this single ingredient we enable a much larger network to train stably at
high replay ratios, yielding super-human Atari-100k median performance.''
\end{quote}
\textbf{Commentary.} BBF is the cleanest evidence that the next decade of Deep RL
improvements will come from \emph{deep-learning hygiene} (normalisation, residuals,
regularisation, resets) rather than from new algorithmic ideas. The companion work
SimBa (Lee 2024) makes the same argument for continuous control. Both papers are
prerequisites for Week 10's modern actor-critic best-practice checklist.

\bigskip
\textit{Sourced from:} \texttt{research/deeprl\_core\_sota.md} \S1.2--1.12 (DQN family
canon and BBF descendants); the OPE-PER connection is from \S0 of
\texttt{research/bandits\_matching\_sota.md}.

\end{document}
